% Part: lambda-calculus
% Chapter: introduction
% Section: overview

\documentclass[../../../include/open-logic-section]{subfiles}

\begin{document}

\olfileid{lam}{int}{ovr}
\olsection{ઝાંખી}

Alonzo Churchએ 1930ના દાયકાની શરૂઆતમાં લૅમ્બડા કલનની રચના મૂળે
રચનાત્મક તર્કશાસ્ત્રના પાયા તરીકે કરી હતી; સંગણનીય વિધેયોના
નિદર્શન તરીકે \emph{નહીં}. પરંતુ ટૂંક સમયમાં બતાવવામાં આવ્યું કે
તે સંગણનીયતાની બીજી વ્યાખ્યાઓ, જેમ કે ટ્યુરિંગ-સંગણનીય વિધેયો
અને આંશિક પુનરાવર્તી વિધેયો, સાથે સમકક્ષ છે. શરૂઆતમાં આ વાતથી
થોડું આશ્ચર્ય થયું હતું, અને તેથી આ લાક્ષણિકીકરણ વધુ રસપ્રદ બને છે.

લૅમ્બડા સંકેતન વિધેયને નામ આપ્યા વિના, તેને વ્યાખ્યાયિત કરતી
સંકેતી અભિવ્યક્તિ દ્વારા સીધો ઉલ્લેખ કરવાની અનુકૂળ રીત છે.
``$f(x) = x + 3$થી વ્યાખ્યાયિત વિધેય~$f$ લો'' એમ કહેવાને બદલે
``$f$ એ વિધેય~$\lambd[x][(x + 3)]$ છે'' એમ કહી શકાય. બીજા શબ્દોમાં,
$\lambd[x][(x+3)]$ એ તેની દલીલમાં ત્રણ ઉમેરતા વિધેયનું માત્ર
\emph{નામ} છે. આ અભિવ્યક્તિમાં~$x$ અસ્થાયી ચલ અથવા સ્થાનધારક છે:
એ જ વિધેયને $\lambd[y][(y + 3)]$ વડે પણ દર્શાવી શકાય છે. બીજા
પ્રાચલો હોય ત્યારે પણ આ સંકેતન ચાલે છે. દાખલા તરીકે,
$g(x, y)$ બે ચલોનું વિધેય હોય અને~$k$ પ્રાકૃતિક સંખ્યા હોય,
તો~$\lambd[x][g(x,k)]$ એ દરેક~$x$ને~$g(x, k)$માં મોકલતું વિધેય છે.

સંકેતી અભિવ્યક્તિ પરથી વિધેય વ્યાખ્યાયિત કરવાની આ રીત
\emph{લૅમ્બડા અમૂર્તીકરણ} કહેવાય છે. તેનું બીજું પાસું
\emph{અનુપ્રયોગ} છે: વિધેય~$f$ (માનો કે પ્રાકૃતિક સંખ્યાઓ પર
વ્યાખ્યાયિત) હોય તો તેને 2 જેવા મૂલ્ય પર \emph{લાગુ} કરી શકાય.
પરંપરાગત સંકેતનમાં તેના પરિણામ માટે આપણે~$f(2)$ લખીએ છીએ.

લૅમ્બડા અમૂર્તીકરણ અને અનુપ્રયોગ જોડીએ તો શું થાય? અમૂર્ત કરેલા
ચલની જગ્યાએ લાગુ કરાયેલી દલીલ મૂકીને મળેલી અભિવ્યક્તિને સરળ બનાવી
શકાય. દાખલા તરીકે,
\[
(\lambd[x][(x + 3)])(2)
\]
ને~$2 + 3$ તરીકે સરળ બનાવી શકાય.

અહીં સુધી આપણે પરંપરાગત ખ્યાલો માટે નવાં સંકેતનો જ આપ્યાં છે.
પરંતુ લૅમ્બડા કલન ગણસિદ્ધાંતના દૃષ્ટિકોણથી વધુ મૂળભૂત ફેરફાર
કરે છે. આ માળખામાં:
\begin{enumerate}
\item દરેક વસ્તુ વિધેય દર્શાવે છે.
\item લૅમ્બડા અમૂર્તીકરણથી વિધેયો વ્યાખ્યાયિત કરી શકાય છે.
\item કોઈપણ વસ્તુને બીજી કોઈપણ વસ્તુ પર લાગુ કરી શકાય છે.
\end{enumerate}
દાખલા તરીકે, લૅમ્બડા કલનમાં~$F$ પદ હોય તો~$F(F)$ હંમેશાં અર્થપૂર્ણ
માનવામાં આવે છે. આ ઉદાર માળખું \emph{પ્રકારવિહીન} લૅમ્બડા કલન
કહેવાય છે. અહીં ``પ્રકારવિહીન''નો અર્થ છે કે કઈ વસ્તુને કઈ વસ્તુ
પર લાગુ કરી શકાય તેના પર કોઈ નિયંત્રણ નથી.

\begin{digress}
લૅમ્બડા કલનનું \emph{પ્રકારિત} સ્વરૂપ પણ છે; તે પ્રકારવિહીન
સ્વરૂપનો મહત્વનો પ્રકારાંતર છે. ઘણી બાબતોમાં પ્રકારિત લૅમ્બડા કલન
પ્રકારવિહીન સ્વરૂપ જેવું હોવા છતાં, તેને પરંપરાગત ગણસિદ્ધાંતના
માળખા સાથે સુસંગત કરવું ઘણું સહેલું છે અને તેના કેટલાક ગુણધર્મો
તદ્દન જુદા છે.

લૅમ્બડા કલન વિશેનું સંશોધન સૈદ્ધાંતિક કમ્પ્યુટર વિજ્ઞાન અને
પ્રોગ્રામિંગ ભાષાઓની રચનામાં કેન્દ્રસ્થાને રહ્યું છે. 1950ના
દાયકામાં John McCarthyએ રચેલી LISP એ આ વિચારોની અસર ધરાવતી
ભાષાનું પ્રારંભિક ઉદાહરણ છે.
\end{digress}

\end{document}
